\section*{Foreword}
\addcontentsline{toc}{section}{Foreword}

Causality is a broad topic, and these lecture notes cover only part of it. 
They originated over a period of three years as a by-product of the course on causality we taught for MSc.\ mathematics students. 
Our aim was to give a mathematically rigorous exposition of the graphical account to causal modeling, reasoning and inference, in the spirit of Wright, Spirtes, Glymour, Scheines, Pearl, and many others. 
Since there seemed to be no book or lecture notes out there that would fit our purpose, we decided to write our own.

The amount of material has grown over the years, and is still growing.
We treat causal modeling with causal Bayesian networks (also known as `DAGs') and structural causal models. 
Some unique features of our exposition are:
\begin{enumerate}
\item we have extended the standard formalisms with \emph{input nodes} to enable a measure-theoretically rigorous treatment of the families of probability distributions that result from perfect interventions;
\item we allow for (sufficiently weak) cycles in structural causal models.
\end{enumerate}
Our treatment is self-contained: we start with the basic definitions (with as prerequisites only basic measure theory and probability theory), and derive everything that is necessary to prove the validity of Markov properties, the do-calculus, adjustment criteria, all the way up to extended versions of the ID algorithm and the FCI algorithm.
We show how---with relatively little extra work---the framework of causal modeling with directed acyclic graphs can be extended to directed graphs that may have cycles.

While the advantages of mathematical rigor should be obvious, the price paid is that the non-trivial conceptual issues are sometimes clouded by technicalities.
We believe that our treatment fills a much needed gap in the literature on causality, and consider it complementary to the many existing writings on similar topics (which often focus more on concepts and less on mathematical rigor).

We are indebted to our teaching assistants Leon Lang, Philip Boeken, Pim de Haan and Noud de Kroon for providing feedback and for spotting several errors in earlier drafts.
While the current version undoubtedly still contains mistakes, we believe that it is now ready for wider exposure. 
We appreciate any feedback that the reader may have, be it on content, typos, or (we hope not) more serious mistakes.


\bigskip
\bigskip
{\flushright Joris Mooij \& Patrick Forr\'e\\
Amsterdam\\
June 2023\\
}
